\chapter{Analyse et conception}

\section{Introduction}
Ce chapitre formalise les besoins du projet puis presente la conception fonctionnelle et technique de \ProjectName. L'objectif est de traduire la problematique metier en un systeme coherent, evolutif et securise.

\section{Specification des besoins}
\subsection{Acteurs}
Deux acteurs principaux sont identifies :
\begin{itemize}[leftmargin=1.2cm]
	\item \textbf{Administrateur technique} : gere le catalogue, les URLs, les utilisateurs et la supervision operationnelle.
	\item \textbf{Analyste metier} : consulte les indicateurs, compare les prix et interprete les tendances.
\end{itemize}

\subsection{Besoins fonctionnels}
Les besoins fonctionnels retenus sont :
\begin{itemize}[leftmargin=1.2cm]
	\item extraction automatisee des prix depuis les sites cibles ;
	\item transformation et nettoyage des donnees (ETL) ;
	\item historisation hebdomadaire et normalisation en EUR ;
	\item visualisation interactive des prix et tendances ;
	\item comparaison par pays, enseigne et magasin ;
	\item administration complete (produits, URLs, utilisateurs) ;
	\item suivi de la qualite de collecte (Operational Visibility).
\end{itemize}

\subsection{Besoins non fonctionnels}
\begin{itemize}[leftmargin=1.2cm]
	\item \textbf{Fiabilite} : controles de validite sur les donnees entrantes.
	\item \textbf{Securite} : authentification et controle des acces par role.
	\item \textbf{Performance} : temps de reponse acceptable pour un usage quotidien.
	\item \textbf{Scalabilite} : possibilite d'ajouter de nouveaux pays et nouveaux sites.
	\item \textbf{Maintenabilite} : architecture modulaire et claire.
\end{itemize}

\section{Analyse conceptuelle (UML)}
\subsection{Diagramme de cas d'utilisation}
Le diagramme de cas d'utilisation met en evidence les interactions entre les deux acteurs et les fonctions principales de la plateforme.

\figureplaceholder{Diagramme UML de cas d'utilisation (Admin technique / Analyste metier)}{Cas d'utilisation principaux de \ProjectName}

\subsection{Diagramme de sequence}
Le diagramme de sequence principal represente le flux de donnees de bout en bout : source web -> scraping -> ETL -> base PostgreSQL -> API FastAPI -> interface React.

\figureplaceholder{Diagramme UML de sequence du flux de donnees}{Sequence globale de traitement et de consultation}

\section{Conception de la base de donnees}
\subsection{Dictionnaire de donnees (extrait)}
\begin{table}[H]
\centering
\caption{Extrait du dictionnaire de donnees}
\begin{tabular}{p{3.5cm}p{4.0cm}p{5.9cm}}
\toprule
\textbf{Entite} & \textbf{Attribut cle} & \textbf{Description} \\
\midrule
product\_formats & product\_format\_id & variante produit (brand/category/range/format/packaging) \\
product\_urls & url, is\_active & source de scraping associee a un produit/site/magasin \\
raw\_staging & payload, status, scraped\_at & donnees brutes collecte + etat de traitement \\
weekly\_price\_summary & week\_start, avg\_price, currency & consolidation hebdomadaire exploitee par le dashboard \\
\bottomrule
\end{tabular}
\end{table}

\subsection{Modele logique de donnees (MLD)}
La conception adopte une separation claire entre :
\begin{itemize}[leftmargin=1.2cm]
	\item donnees de reference (catalogue, sites, magasins) ;
	\item donnees brutes de collecte ;
	\item donnees nettoyees et analytiques.
\end{itemize}

\figureplaceholder{Modele logique de donnees (MLD) du projet}{Schema logique de la base PostgreSQL}

\subsection{Regles de conception}
\begin{itemize}[leftmargin=1.2cm]
	\item cles etrangeres pour garantir la coherence relationnelle ;
	\item contraintes d'unicite sur les combinaisons critiques ;
	\item index sur les colonnes de recherche frequente ;
	\item regles de qualite (exclusion des prix <= 0).
\end{itemize}

\section{Architecture logique de la solution}
La solution est concue en couches : collecte, ETL, stockage, API, presentation. Cette decomposition permet d'isoler les responsabilites et de simplifier la maintenance.

\figureplaceholder{Architecture logique en couches}{Conception modulaire de \ProjectName}

\section{Conclusion}
L'analyse et la conception realisent le lien entre besoins metier et implementation technique. Elles fournissent une base robuste pour une realisation progressive et maitrisee de la plateforme.

\clearpage
